\subsection{Retail execution: state scope and a visible-call counterexample}\label{sec:tau-cover-bridge}
We replay the 115 public reference programs in the pinned $\tau$-bench retail simulator. There is no LLM planner or user simulator in this experiment. Ninety-seven tasks have an initially pending order owned by the task user; 18 have none and receive no cover. Three seeds select orders, and cover is inserted before or after the prescribed program. Utility uses 97 task clusters, each with three repeated units per position. Policies include no cover, discard, shadow execution, read only, address self-write, item self-write, instruction-aware selection, other-account cancellation and blind cancellation. Destructive conditions are local negative controls, not proposed deployment policies.

\paragraph{Separate the outcome scopes}
Table~\ref{tab:tau-cover-bridge} reports reference database equality, full active-account equality and equality on objects \emph{changed} by the reference program. The last scope does not include every read/accessed object. The database criterion corresponds to the state part of the native evaluator; the extension script does not evaluate generated response correctness and we do not relabel its state equality as an end-to-end native task reward. All-account equality now checks every order owned by the active user and that user's record, rather than only initially pending orders.
\begin{table}[t]\centering\footnotesize
\caption{Retail state and observer results. State columns use 97 eligible tasks (291 units per position); observer columns use 49 held-out tasks (147 units). Percentages are task means across three seeds. B/A denotes before/after.}\label{tab:tau-cover-bridge}
\setlength{\tabcolsep}{2.2pt}
\begin{tabularx}{\columnwidth}{@{}Yrrrr@{}}\toprule
\multicolumn{5}{@{}l}{\textbf{A. Distinct state scopes (reference execution as comparator)}}\\
Policy / position & DB eq. & Account eq. & Object eq. & Mutated\\\midrule
No cover B/A & 100.00 & 100.00 & 100.00 & 0.00\\
Discard B/A & 100.00 & 100.00 & 100.00 & 0.00\\
Shadow B/A & 100.00 & 100.00 & 100.00 & 0.00\\
Read only B/A & 100.00 & 100.00 & 100.00 & 0.00\\
Address self-write B/A & 100.00 & 100.00 & 100.00 & 0.00\\
Item self-write B & 40.55 & 40.55 & 70.10 & 59.45\\
Item self-write A & 67.70 & 67.70 & 97.25 & 32.30\\
Instruction aware B & 35.05 & 35.05 & 78.35 & 70.10\\
Instruction aware A & 46.39 & 46.39 & 89.69 & 53.61\\
Other-account cancel B & 0.00 & 100.00 & 100.00 & 100.00\\
Other-account cancel A & 0.00 & 100.00 & 100.00 & 100.00\\
Blind cancel B & 11.34 & 11.34 & 62.54 & 100.00\\
Blind cancel A & 41.24 & 41.24 & 92.44 & 58.76\\
\midrule
\multicolumn{5}{@{}l}{\textbf{B. Legacy learned observer versus allowed-view rules}}\\
Policy / position & $D_0$ acc. & $D_2$ acc. & $D_2$ BA & Rule acc.\\\midrule
Majority (all conditions) & 77.55 & 77.55 & 50.00 & ---\\
Instruction aware B & 51.02 & 73.47 & 47.37 & 100.00\\
Instruction aware A & 51.02 & 71.43 & 52.51 & 100.00\\
Other-account cancel B & 22.45 & 71.43 & 46.05 & 100.00\\
Other-account cancel A & 22.45 & 71.43 & 46.05 & 100.00\\
Blind cancel B & 22.45 & 71.43 & 46.05 & 100.00\\
Blind cancel A & 22.45 & 68.03 & 47.09 & 100.00\\
\bottomrule\end{tabularx}
\par\smallskip\raggedright
DB eq.: equality of orders/users against reference (products are read-only), the native reward's database criterion, not an output-response/authorization check. Account eq.: all orders owned by the active user plus that user's record. Object eq.: objects changed by the reference program, not all accessed objects. Mutated: cover itself changes state. Other-account calls yield 0\% DB equality but 100\% active-account equality. BA: balanced accuracy. Rule: remove the known first/last cover call, then inspect cancellation verbs; blind/other-account cancellation also admits the $N_{\rm cancel}\ge2$ rule. Both rules attain 100\% acc./BA in their stated conditions. Other non-cancellation policies retain 100\% verb inference. Full intervals, AUROC and confusion matrices appear in Supplement S9 and the JSON artifact.
\end{table}


Address self-writes preserve the reference state in this simulator. Item self-writes mutate state in 173/291 pre-insertion units (59.45\%): the tool changes the status to \texttt{pending (item modified)} and appends a refund entry even when the item identifier and price are unchanged. The remaining 118/291 calls fail because the requested variant is unavailable, explaining why this mutation rate is below blind cancellation's 291/291 successful calls. This is a state-preservation counterexample from one API, not a prevalence estimate or a discovery about idempotence. Formally, a no-change requirement $f(s,a)=s$ differs from repeated-call idempotence $f(f(s,a),a)=f(s,a)$.

Instruction-aware cover uses only instruction text and the initial order table for selection, with address self-write as fallback. It preserves reference-changed objects in 78.35\%/89.69\% of before/after units, but preserves the full account and database in only 35.05\%/46.39\%. Cancelling an unmentioned order can therefore damage a bystander resource. The heuristic was developed using the available task collection; it has no held-out policy-development evaluation. Its 68/97 availability (70.10\%) and the item-self-write object equality of 204/291 (70.10\%) are separate counts that happen to yield the same percentage.

The other-account condition selects a pending order belonging to another simulated user. It is not a provisioned dedicated decoy account. It preserves the active user's complete account and reference-changed objects in these tasks, while changing the shared database in every eligible unit (0\% database equality). Thus account isolation of an effect cannot be described as whole-database preservation or an authorized production defense.

\paragraph{Observer permissions and simple attacks}
The secret label is whether the reference program \emph{attempts} a cancellation, including a call that may fail after cover execution. The observer sees the ordered dispatched tool names, write counts and response lengths; failed calls remain visible. It sees no arguments, user/order identifiers or database state. We do not infer that a passive TLS observer has these permissions. A backend with plaintext state would instead need a different confidentiality target.

The fixed task split contains 48 training tasks (7 positive, 41 negative) and 49 test tasks (11 positive, 38 negative); seeds stay within tasks. Always predicting no cancellation gives 77.55\% test accuracy and 50\% balanced accuracy. The legacy $D_2$ linear classifier uses cancellation presence, write count and write-response bytes. Its 68.03--73.47\% accuracy under cancellation cover is below the majority baseline, and its balanced accuracy ranges from 45.61\% to 52.79\%. An increase over $D_0$ is therefore insufficient evidence for its claimed privacy interpretation.

For blind and other-account cancellation, exactly one extra cancellation is dispatched in every eligible session. The permitted rule $\widehat y=\mathbf{1}[N_{\rm cancel}\ge2]$ recovers the reference label at 100\% accuracy and balanced accuracy on all 49 test tasks and three seeds, before and after insertion. For any policy with one live cover call at a known first/last position, removing that call and testing for a remaining cancellation also gives 100\%. It needs no identifiers, no response content and no training labels. These are measured finite-sample results and a structural property of the fixed insertion design, not proof of general attack performance. Identifier filtering is unnecessary in this view and is not claimed as an evaluated attack.

\paragraph{Terminal equality does not establish path properties}
Blind pre-insertion cancellation leaves the final database equal to reference in 33/291 units across 12 tasks, even though it preempts the later reference cancellation. The exact per-task counts are in Supplement S9. After insertion, 120/291 cover calls are rejected because reference execution already changed the order's status, and those rejections explain the 41.24\% equality rate. Rejection is not successful benign cover. In the 82 tasks with no reference execution error (246 units), before/after equality is 24/246 (9.76\%) and 94/246 (38.21\%). These results illustrate the limits of endpoint scoring; consent and authorization were not measured.

\soktakeaway{SoK Takeaway 4 (State preservation and inference are separate)}{Same-value tool inputs can change workflow state. Equality on reference-changed objects can miss damage elsewhere in an account; account equality can coexist with global database changes. Even a state-preserving route need not hide a cancellation label from the specified dispatch observer. The tested policies do not establish a privacy--utility frontier.}
